\documentclass[11pt]{article}
\usepackage[margin=1in]{geometry}
\usepackage{amsmath,amssymb}
\usepackage{enumitem}
\usepackage{booktabs}

\newcommand{\indep}{{\bot\negthickspace\negthickspace\bot}}
\DeclareMathOperator{\E}{\mathbb{E}}

\pagestyle{plain}

\begin{document}

\begin{center}
{\Large \textbf{PS200C: Quantitative Methods}} \\[6pt]
{\large In-class Activity: Selection on Observables}
\end{center}

\vspace{4pt}
\noindent\textbf{Name(s):} \hrulefill

\vspace{14pt}

%%% QUESTION 1 %%%
\noindent\textbf{Question 1: The identification argument}

\medskip
\noindent Suppose you assume conditional ignorability: $\{Y_{1i}, Y_{0i}\} \indep D_i \mid X_i$.

\begin{enumerate}[label=(\alph*)]
\item In one sentence, what does this assumption say in plain English?

\vspace{2cm}

\item What are some other names for this assumption?

\vspace{2cm}
\end{enumerate}


%%% QUESTION 2 %%%
\noindent\textbf{Question 2: Evaluating SOO claims}

\medskip
\noindent For each study, state what the SOO/CI assumption would require, and whether you find it plausible.

\begin{enumerate}[label=(\alph*)]
\item A researcher estimates the effect of college attendance on earnings by comparing college-goers to non-goers, ``controlling for'' high school GPA, parental income, and state of residence.

\vspace{2.5cm}

\item Blattman \& Annan (2009) estimate the effect of abduction by the LRA on education, conditioning on village location and age, arguing that abduction was quasi-random conditional on these variables.

\vspace{2.5cm}
\end{enumerate}

\newpage

%%% QUESTION 3 %%%
\noindent\textbf{Question 3: Post-treatment variables}

\medskip
\noindent A randomized job training program ($D$) aims to increase earnings ($Y$).  A researcher notices that some treated individuals dropped out of the program early, and restricts the analysis to ``completers'' vs.\ controls.

\begin{enumerate}[label=(\alph*)]
\item Is ``program completion'' a pre-treatment or post-treatment variable?

\vspace{1cm}

\item Why could restricting to completers bias the estimate, even though $D$ was randomly assigned?  (Think about what kinds of people complete vs.\ drop out.)

\vspace{2.5cm}

\end{enumerate}


%%% QUESTION 4 %%%
\noindent\textbf{Question 4: Subclassification}

\medskip
\noindent You have two strata defined by a covariate $X \in \{\text{low}, \text{high}\}$:

\begin{center}
\begin{tabular}{@{} l c c c @{}}
\toprule
Stratum & $\bar Y$ (treated) & $\bar Y$ (control) & $N_j$ (total) \\
\midrule
$X = \text{low}$  & 12 & 8 & 60 \\
$X = \text{high}$ & 20 & 18 & 40 \\
\midrule
Total & & & 100 \\
\bottomrule
\end{tabular}
\end{center}

\begin{enumerate}[label=(\alph*)]
\item Compute the subclassification estimate of the ATE.

\vspace{2cm}

\item Suppose the treated group is disproportionately in the ``high'' stratum.  Would you expect the unadjusted DIM to overestimate or underestimate the ATE?  Why?

\vspace{2.5cm}

\item Under what assumption is your answer to (a) a valid causal effect?

\vspace{2cm}
\end{enumerate}

\end{document}
